\section{Nonlinear corrections to a complete affine baseline}\label{sec:residuals}

An operator bound becomes useful for a computation only after its
\emph{complete} baseline is exhibited. This appendix gives that conditional
interface. It allows arbitrary nonlinear corrections and arbitrary correctness
sets, but supplies no cheap affine approximation to an arbitrary machine.
Fixed-advice fibre counting is already explicit in
\cite[Theorem~5.5]{ClassUniformMemo2026} and
\cite[Corollary~6.16]{SeventeenthPaper2026}. We use the same counting principle
with a chronological computation and retain the actual propagated correction
span. The distinction prevents charging one correction again at each consumer.

\subsection{The complete baseline and its joint record}

Let $F$ be a finite field of size $q$, and let $X\in F^N$ be the raw input.
All current payload, retained control and earlier output coordinates belong
to the state. A finite sequence of state spaces can be embedded in their
direct sum, so we use one finite-dimensional state space $V$. The baseline
initialization, updates and readout are global affine maps
\[
 v_0=JX+j,\qquad G_i(v)=L_i v+g_i,\qquad \operatorname{out}=Ov_m+o.
\]
Let the complete baseline output be
$K_0X+c_0$. At a marked update the actual map may be arbitrary. Its residual
on the actual old state is $E_i u_i(X)$, where $E_i$ is a fixed linear
insertion map. At other updates the residual is zero. The initialization must
also agree with the baseline, or its departure must be included as a marked
initial residual. All unmarked maps, including final readout, remain the same
\emph{global} affine maps on every reached state.

Choose one preimage $u_i(X)$ in a fixed linear section of $\operatorname{im}E_i$.
These values need not be independent or linear in $X$. Write $Z(X)$ for their
one joint tuple, and put
\[
 A=\tau-K_0,\qquad D=\dim\ker A,\qquad
 W=\sum_i\operatorname{im}(OL_{m-1}\cdots L_{i+1}E_i),
\]
where the empty product is the identity. These are comparison records and
subspaces; the definition does not give a native machine free access to them.

\begin{theorem}[Joint residual replay]\label{thm:residual-replay}
\formal{ResidualPropagation.lean}{residual_readout_sum,actual_readout_sum,actual_readout_insertions}
The complete-program counting endpoint is
\lean{actual_correct_card_le} in \texttt{ResidualProgram.lean}.
For every fixed tuple $z=(u_i)$, replay gives
\[
 \operatorname{out}_z(X)=K_0X+c_0+
       \sum_i OL_{m-1}\cdots L_{i+1}E_i u_i.
\]
On an arbitrary correctness set $C$, the actual tuple $Z(X)$ reproduces the
actual computation. Consequently
\begin{equation}\label{eq:residual-record}
 |C|\le |Z(C)|q^D.
\end{equation}
\end{theorem}

\begin{outline}
Induct in the original update order. After fixing the residual values, each
update is $v\mapsto L_iv+g_i+E_i u_i$, so its linear part is still $L_i$.
Expanding the recurrence propagates each inserted residual through exactly
the subsequent linear maps and then through $O$. With $z=Z(X)$ the
recurrence agrees with every actual update on $X$, including retained
outputs. On a fixed record fibre, correctness is the affine equation
$AX=c(z)$. Subtract one solution, if there is one, to inject the fibre into
$\ker A$. It therefore contains at most $q^D$ elements. Sum over the
\emph{attained} joint records. The formal proofs perform this recurrence and
translate-to-kernel injection; they do not assume that $C$ is a subspace.
\end{outline}

\begin{theorem}[Joint output-span bound]\label{thm:residual-span}
\formal{ResidualSpan.lean}{residualSpanDescriptor_injective,finrank_preimage_le_kernel_add_intersection,residual_span_card_le}
\formal{ResidualOutputSpan.lean}{actual_correct_span_card_le}
For the actual program, and more generally, if $C\subseteq F^N$ satisfies $AX-c\in W$ for every $X\in C$,
then
\begin{equation}\label{eq:residual-span}
 |C|\le q^{D+\dim(W\cap\operatorname{im}A)}
       \le q^{D+\dim W}.
\end{equation}
For the replay above, take $c=c_0$ and the displayed propagated span $W$.
\end{theorem}

\begin{outline}
Subtract one point of $C$ to inject it into $A^{-1}(W)$. Choose a linear
projection onto $\ker A$ that is the identity on that kernel. The map sending
$x\in A^{-1}(W)$ to its projection and $Ax$ is injective, with values in
$\ker A\times(W\cap\operatorname{im}A)$. This proves the dimension bound
and then the finite-field cardinality bound. In particular, the invoice
uses the joint propagated span, not a sum that charges dependent copies
repeatedly. Its dimension is at most the sum of the insertion ranks.
\end{outline}

\begin{corollary}[Information and sparse corrections]\label{cor:residual-information}
\formal{ResidualEntropy.lean}{residual_entropy_lower,residual_uniform_entropy_lower,actual_correct_entropy_lower,residual_span_information}
The sparse-count endpoint is \lean{actual_sparse_card_le} in
\texttt{ResidualSparse.lean}.
For any law supported on $C$, with entropy in base $q$,
\[
 H_q(Z)\ge H_q(X)-D.
\]
If $C$ is uniform and $\Delta=N-\log_q|C|$, this gives
$H_q(Z)\ge N-D-\Delta$ and $\dim W\ge N-D-\Delta$.
If each of $m$ public insertion ports has rank at most $r$, the correctly
handled inputs on which at most $s$ residuals are nonzero number at most
\begin{equation}\label{eq:sparse-residual-menu}
 q^D\sum_{j=0}^{\min(s,m)}\binom mj(q^r-1)^j.
\end{equation}
\end{corollary}

\begin{outline}
Each conditional input fibre has at most $q^D$ points, so
$H_q(X\mid Z)\le D$. Since $Z$ is a function of $X$, the entropy chain rule
gives the first inequality. The span bound gives the second assertion for
uniform $C$. For the sparse count, choose coordinates in each insertion
image, with the zero residual represented uniquely by zero. A support of
$j$ positions allows at most $(q^r-1)^j$ nonzero value tuples. Sum over
supports and apply \cref{eq:residual-record}. The locations may depend
arbitrarily on $X$; their support count is already included. The formal
count records each actual residual as a value in its insertion image,
so it charges the insertion rank and does not assume a particular
preimage selection. The shared support/value encoding also permits
different alphabets at different ports.
\end{outline}

For fixed $q,r$, $m=O(N)$ and $s=o(N)$, the logarithm of the sum in
\cref{eq:sparse-residual-menu} is $o(N)$. The formal endpoint is
\lean{LeanMachine.sparseMenu_log_div_tendsto_zero}. Its finite proof multiplies
the sum by $x^s$ and bounds the result by $(1+Ax)^m$ for $0<x\le1$,
where $A=q^r-1$.
Choose $x$ small enough that $m\log(1+Ax)$ costs at most half the desired
linear budget, then bound $s\log(1/x)$ by the other half. This treats
$s=0$ directly and requires no lower bound on $m/N$.

\subsection{What can be substituted, and what remains missing}

Suppose $N=bd^2$ and the \emph{complete} baseline has its column factors in
one common triangularizable algebra. The flag bound and its entry-width
extension give $D\le bd(d+1)/2$\cite[corollary ``One triangularizable side''
and remark ``Entries of $b$ bits'']{SeventeenthPaper2026}.
Thus correct inputs of deficit $\Delta$ require
\[
 H_q(Z)\ge \frac{bd(d-1)}2-\Delta,\qquad
 \dim W\ge \frac{bd(d-1)}2-\Delta.
\]
An independently established sharper operator bound may replace this $D$.
With fixed-width insertion events, sublinear nonzero corrections therefore
cover at most $q^{(1/2+o(1))N}$ inputs. A whole-row correction of rank $O(db)$
is not a fixed-width event.

The revised entry-rule bound of
\cite[Proposition~6.15 and Corollary~6.16]{SeventeenthPaper2026} is another
available specialization when its actual entry functionals are supplied.
It prices advice by its support and changes of an entry decoder by half a
symbol. A change of a whole row is $d$ entry changes. Additive advice alone
need not change any decoder. Our rank invoice and that change count are
different statistics; neither may be substituted for the other without
checking the complete rule.

An input-dependent branch or schema can be included in the joint record.
Then each branch has its own complete baseline and kernel bound $D_\theta$;
the support estimate sums the corresponding $q^{D_\theta}$ fibre bounds,
and entropy gives $H_q(Z)\ge H_q(X)-\mathbb E D_\theta$.
These attained-branch and entropy bounds are
\lean{residual_branch_card_le} and
\lean{residual_branch_entropy_lower} in \texttt{ResidualBranches.lean}. Unknown earlier
encodings, address histories, mutable flags or nonlinear unmarked readout
must be included in such a proved interface. They cannot be omitted by
checking only the zero-residual trajectory.

These statements exclude sparse corrections to an \emph{exhibited}
structured computation. They do not show that every fast transposer is
close to one, or that producing a correction of large rank costs superlinear
time. Supplying that structural comparison with a small joint record
remains the construction/impossibility bridge.

\subsection{Two sharp complete-baseline examples}

\begin{proposition}[Separable baseline defect]\label{prop:separable-defect}
\formal{ResidualSeparable.lean}{Separable.rank_defect_lower}
\formal{ResidualSeparableSharp.lean}{Separable.initialProjection_rank,Separable.rank_defect_sharp}
For $L(X)=AXB^{\mathsf T}$ over any field, put
$r=\min(\rank A,\rank B)$. Then
\[
 \rank(\tau-L)\ge d^2-\frac{r(r+1)}2.
\]
Equality is attained for $A=B=\operatorname{diag}(I_r,0)$ in every
characteristic. Thus this is a sharp bound for one separable baseline,
not for a sum of arbitrary tensor terms.
\end{proposition}
\begin{outline}
Writing $Z=X^{\mathsf T}$ turns agreement into $Z=AZ^{\mathsf T}B^{\mathsf T}$.
If $a=\rank A$, choose invertible $U,V$ with $UAV^{-1}=E_a$.
For $Y=UZV^{\mathsf T}$ and $D=VBU^{-1}$ the equation is
$Y=E_aY^{\mathsf T}D^{\mathsf T}$. The last $d-a$ rows vanish, and the
upper-left $a$-square block $Q$ determines the others and satisfies
$Q=Q^{\mathsf T}D_{11}^{\mathsf T}$. Its lower triangle, including the
diagonal, determines $Q$: it gives column zero and therefore row zero;
the recovered rows and the lower triangle then give the next column and
row, successively. The formal proof obtains the invertible factors and
the active coordinate reindexing from the actual matrix rank normal form.
The resulting restriction is a linear injection into the unordered-pair
coordinates of the active block, so the kernel has dimension at most
$a(a+1)/2$. Transposing the
agreement equation exchanges $A,B$ and gives the same bound with
$a=\rank B$. Rank--nullity proves the claim. For $A=B=E_r$, agreement
is exactly symmetry inside the leading $r$-square block and zero outside,
with $r(r+1)/2$ independent entries.
\end{outline}

\begin{proposition}[Nonlinear row-then-column agreement]\label{prop:row-column-agreement}
\formal{ResidualRowColumn.lean}{RowColumn.agreement_card_le,RowColumn.helper_agreement_card_le,RowColumn.identity_agreement_card}
\formal{ResidualRowColumnSharp.lean}{RowColumn.exists_sharp_helper,RowColumn.helper_bound_attained}
Over an alphabet of size $q\ge2$, first apply an arbitrary fixed map to
each row, then an arbitrary fixed map to each column. The resulting
complete map agrees with transpose on at most $q^{d(d+1)/2}$ inputs.
If one joint helper chooses all these maps, at most
$|\operatorname{range}H|q^{d(d+1)/2}$ inputs can be correct.
Both bounds are sharp as abstract counting statements.
\end{proposition}
\begin{outline}
Write the intermediate matrix as $Y_{j,:}=R_j(X_{j,:})$ and the output
column as $C_j(Y_{:,j})$. On agreement inputs,
$X_{j,:}=C_j(Y_{:,j})$ and therefore
$Y_{j,:}=R_j(C_j(Y_{:,j}))$. The lower triangle of $Y$ determines column
zero, hence row zero; then column one, hence row one; continue in order.
It determines all of $Y$, which in turn determines $X$. This injection
proves the first count without assuming either family is linear or
injective. Fix the helper and sum to prove the second count. Identity
row/column maps attain the first bound on symmetric inputs. For the
second, identify the alphabet with the cyclic group of order $q$ and
let the helper give $X_{ij}-X_{ji}$ for $i<j$. Identity row maps followed
by column maps adding or subtracting the corresponding differences
transpose every input. Every helper tuple is attained and there are
$q^{d(d-1)/2}$ tuples. Computing and delivering this helper is not free
in a native model.
\end{outline}
